\documentclass{article}
\usepackage[UTF8]{ctex}
\usepackage{amsmath,amssymb,amsfonts}
\usepackage{enumitem}
\usepackage{geometry}
\geometry{a4paper, margin=1in}
% === 错题本排版优化（不参与题目标识） ===
\setlength{\parindent}{0pt}
\setlength{\parskip}{0.35em plus 0.10em minus 0.08em}
\setlength{\emergencystretch}{3em}
\setlist{itemsep=0.25em, topsep=0.40em, parsep=0pt}
\allowdisplaybreaks[2]
% === 错题本排版优化结束 ===
\begin{document}
\section*{郑州大学2024年高等代数真题}
\section*{118．郑州大学2024年高等代数试卷}
一．填空题．每题5分，共30分．\\

\subsection{题目1}
1．\(\left|\begin{array}{cccc}1 & -1 & 1 & x-1 \\ 1 & -1 & x+1 & -1 \\ 1 & x-1 & 1 & -1 \\ x+1 & -1 & 1 & -1\end{array}\right|=\)

\subsection{题目2}
2．设 \(n\) 阶矩阵 \(A=\left(\begin{array}{ccccc}1 & a & a & \cdots & a \\ a & 1 & a & \cdots & a \\ \vdots & \vdots & \vdots & & \vdots \\ a & a & a & \cdots & 1\end{array}\right), n \geq 3, A\) 的秩为 \(n-1\) ，则 \(a=\) \(\_\_\_\_\)

\subsection{题目3}
3．设 \(A=\left(\alpha_{1}, \alpha_{2}, \alpha_{3}, \alpha_{4}\right)\) 为 4 阶矩阵，\(A^{*}\) 为 \(A\) 的伴随矩阵，若单个向量 \(\beta \neq 0\) 是方程组 \(A X=0\) 的一个基础解系，则 \(A^{*} X=0\) 的基础解系含有解向量的个数是 \(\_\_\_\_\)。

\subsection{题目4}
4．实二次型 \(f\left(x_{1}, x_{2}, x_{3}\right)=\left(x_{1}+2 x_{2}\right)^{2}+\left(2 x_{2}-x_{3}\right)^{2}+\left(x_{1}+x_{3}\right)^{2}\) 的正惯性指数为 \(\_\_\_\_\)．

\subsection{题目5}
5．设 \(\alpha_{1}=\left(\begin{array}{l}a \\ 1 \\ 1\end{array}\right), \alpha_{2}=\left(\begin{array}{l}1 \\ a \\ 1\end{array}\right), \alpha_{3}=\left(\begin{array}{l}1 \\ 1 \\ a\end{array}\right), \alpha_{4}=\left(\begin{array}{c}1 \\ 1 \\ -2\end{array}\right)\) ，且 \(\alpha_{1}, \alpha_{2}, \alpha_{3}\) 与 \(\alpha_{2}, \alpha_{3}, \alpha_{4}\) 等价，则\(a\) 满足 \(\_\_\_\_\)．

\subsection{题目6}
6．设 \(\sigma\) 是 \(n\) 维线性空间 \(V\) 上的线性变换，\(\sigma\) 有 \(n\) 个不同的特征值 \(\lambda_{1}, \lambda_{2}, \cdots, \lambda_{n}\) ，则 \(V\) 的所有 \(\sigma\) 不变子空间的个数是 \(\_\_\_\_\)．

二．计算与证明．每题15分，共120分．\\

\subsection{题目7}
1．如果实系数多项式 \(f(x)\) 满足 \(f\left(2 x^{2}+1\right)=2 f^{2}(x)+1, \forall x \in \mathbb{R}, f(0)=0\) ，证明：\(f(x) \equiv x\) ．

\subsection{题目8}
2．设 \(A=\left(\begin{array}{ll}1 & a \\ 1 & 0\end{array}\right), B=\left(\begin{array}{ll}0 & 1 \\ 1 & b\end{array}\right)\) ，当 \(a, b\) 为何值时，存在矩阵 \(C\) ，使得 \(A C-C A=B\) ，并求所有矩阵 \(C\) 。

\subsection{题目9}
3．设 \(A\) 是 \(n\) 阶矩阵，\(n>1\) 。如果对任意 \(n\) 阶矩阵 \(B\) 都有 \(|A+B|=|A|+|B|\) ，证明：\(A=O\) 。

\subsection{题目10}
4．已知实二次型 \(f\left(x_{1}, x_{2}, x_{3}\right)=a x_{1}^{2}+a x_{2}^{2}+a x_{3}^{2}+2 x_{1} x_{2}+2 x_{1} x_{3}+2 x_{2} x_{3}\) 用正交替换 \(X=T Y\) 化为标准形 \(f\left(x_{1}, x_{2}, x_{3}\right)=b y_{1}^{2}+c y_{2}^{2}, b, c \neq 0\) ，求 \(a, b, c\) 并写出正交替换及所化成的标准二次型。

\subsection{题目11}
5．求复矩阵 \(A=\left(\begin{array}{ccc}3 & 2 & 1 \\ 1 & 4 & 1 \\ -3 & -6 & -1\end{array}\right)\) 的约当标准形 \(J\) ，并求可逆矩阵 \(T\) ，使得 \(T^{-1} A T=J\) 。

\subsection{题目12}
6．设

\[
\alpha_{1}=\left(\begin{array}{c}
a_{1}+b \\
a_{1} \\
\vdots \\
a_{1}
\end{array}\right), \alpha_{2}=\left(\begin{array}{c}
a_{2} \\
a_{2}+b \\
\vdots \\
a_{2}
\end{array}\right), \cdots, \alpha_{n}=\left(\begin{array}{c}
a_{n} \\
a_{n} \\
\vdots \\
a_{n}+b
\end{array}\right)
\]

其中 \(\sum_{i=1}^{n} a_{i} \neq 0\) 。这 \(W=L\left(\alpha_{1}, \alpha_{2}, \cdots, \alpha_{n}\right)\) ，求 \(W\) 的维数与一组基。

\subsection{题目13}
7．设 \(A\) 是 \(n\) 阶实矩阵，\(A^{2}=E\) ，证明：

\begin{itemize}
  \item[（1）] \(r(A+E)+r(A-E)=n\) ．
  \item[（2）] \(A\) 与对角矩阵相似．
  \item[（3）] \(\mathbb{R}^{n}=V_{1} \oplus V_{2}\) ，其中 \(V_{1}=\left\{X \in \mathbb{R}^{n} \mid(A+E) X=0\right\}, V_{2}=\left\{X \in \mathbb{R}^{n} \mid(A-E) X=0\right\}\) 。
\end{itemize}

\subsection{题目14}
8．设 \(V=P^{2 \times 2}\) 是数域 \(P\) 上的线性空间，\(A=\left(\begin{array}{ll}a & b \\ c & d\end{array}\right) \in V\) ，线性变换 \(\sigma: X \rightarrow A X, \forall X \in P^{2 \times 2}\) ．

\begin{itemize}
  \item[（1）] 求线性变换 \(\sigma\) 在基
\[
E_{11}=\left(\begin{array}{ll}
1 & 0 \\
0 & 0
\end{array}\right), E_{12}=\left(\begin{array}{ll}
0 & 1 \\
0 & 0
\end{array}\right), E_{21}=\left(\begin{array}{ll}
0 & 0 \\
1 & 0
\end{array}\right), E_{22}=\left(\begin{array}{ll}
0 & 0 \\
0 & 1
\end{array}\right)
\]
下的矩阵．
  \item[（2）] 如果 \(A\) 相似于对角矩阵，证明线性变换 \(\sigma\) 在 \(V\) 的某组基下的矩阵是对角矩阵。
\end{itemize}
\end{document}
